\begin{exercise}{10.1}{?}
    Welche Matrizen von $\operatorname{GL}(2; \mathbb C)$ liefern Automorphismen
    von $\mathbb P^1 \mathbb C$, die den unendlichen fernen Punkt alias die Gerade
    durch $(1,0)$ festhalten?
\end{exercise}

Sei $A \in \operatorname{GL}(2;\mathbb C)$ ein solcher Automorphismus. Dann 
muss gelten für alle $\lambda \in \mathbb C^\times$:
\[
    A\begin{pmatrix} \lambda \\ 0 \end{pmatrix} =
    \begin{pmatrix}
        a & b \\
        c & d
    \end{pmatrix} \begin{pmatrix} \lambda \\ 0 \end{pmatrix}
    =
    \begin{pmatrix} \lambda a \\ \lambda c \end{pmatrix}
    \in \mathbb C \begin{pmatrix} 1 \\ 0 \end{pmatrix}
    \quad \Longrightarrow \quad c = 0 \quad \text{und} \quad a \neq 0.
\]
Da $A \in \operatorname{GL}(2;\mathbb C)$ gilt
\[
    \det(A) = ad \neq 0 \quad \Longrightarrow \quad a,d \neq 0.
\]
Somit sind die gesuchten Matrizen Elemente der Menge
\[
    \boxed{\Bigl\{
        \begin{pmatrix}
            a & b \\
            0 & d
        \end{pmatrix} \mid a,d \in \mathbb C^\times, b \in \mathbb C
    \Bigr\}.}
\]

\begin{exercise}{10.2}{?}
    Welche Matrizen von $\operatorname{SL}(2; \mathbb R)$ liefern eine Kongruenz
    der oberen Halbebene, die die Halbgerade $\{3+iy \mid y \geq 4\}$ auf die
    Halbgerade $\{iy \mid y \geq 1\}$ abbilden?
\end{exercise}

Sei enstprechende Matrix mit zugehöriger Möbiustransformation
\[
    A=
    \begin{pmatrix}
        a & b\\
        c & d
    \end{pmatrix}
    \in \operatorname{SL}(2;\mathbb R) \qquad \text{und} \qquad
    f_A(z)=\frac{az+b}{cz+d}.
\]
Beide Halbgeraden laufen für \(y\rightarrow\infty\) gegen den unendlich fernen Punkt \(\infty\).
Da Möbiustransformationen den Rand der oberen Halbebene auf sich
abbilden und die erste Halbgerade auf die zweite abgebildet werden soll,
gilt daher \(f_A(\infty)=\infty\).
Für \(f_A(z)\) ist dies genau dann der Fall, wenn \(c=0\).
Zusätzlich gilt \(d = \frac1a\), weil \(\det(A) = ad = 1\) laut Definition.
Somit hat \(A\) die Form
\[
    A=
    \begin{pmatrix}
        a & b\\
        0 & \frac1a
    \end{pmatrix}.
\]

Der endliche Endpunkt \(3+4i\) der ersten Halbgeraden muss auf den
endlichen Endpunkt \(i\) der zweiten Halbgeraden abgebildet werden.
Also gilt \(f_A(3+4i)=i\).
Damit erhalten wir
\[
    f_A(3+4i)=\frac{a(3+4i)+b}{\frac{1}{a}}=3a^2+ab+4ia^2=i.
\]
Untersuchen wir Real- und Imaginäteil folgt
\[
    4ia^2 = i \quad
    \Longleftrightarrow \quad a^2 = \frac{1}{4} \quad
    \Longleftrightarrow \quad a = \frac{1}{2}, \quad \text{und}
\]
\[
    3a^2+ab = 0 \quad 
    \Longleftrightarrow \frac{3}{4}+\frac{b}{2} = 0 \quad
    \Longleftrightarrow b = -\frac{3}{2}.
\]
Nach Satz 2.5.9 sind die gesuchten Matrizen daher genau
\[
    \boxed{
        \pm
        \begin{pmatrix}
            \frac12 & -\frac32\\
            0 & 2
        \end{pmatrix}.
    }
\]
